% ===========================================================================
\chapter{Trust, but verify}\label{ch:checks}

Route planning is the kind of code where a bug hides well. A search that
is wrong now and then still finds a route every time, and the route still
looks reasonable on the map; it is just not quite the best one, for some
pairs of places. Nobody notices until somebody does.

Our contracts already check every single call: every path found is a real
route, its cost adds up, every projection goes back where it came from.
What contracts can't easily say is how different parts relate to each
other, or what should be true of the map as a whole. So the last file of
the core gets a set of \emph{checks}: a program that tests properties we
know must hold, over every pair of places on the island, for both drivers
and walkers.

\section{A checker}

A checker counts what passes and what fails, and prints each result. At
the end it prints a summary, and if anything failed it stops the program
with a non-zero exit code, the signal that scripts and build tools look
for.

\begin{newfile}{checks.nex}
-- checks.nex: executable checks for every piece of the project.
-- Run: nex checks.nex

intern Search
intern Helpers

class Checker
feature
  passed: Integer
  failed: Integer

  check(label: String, ok: Boolean) do
    if ok then
      passed := passed + 1
      say("PASS " + label)
    else
      failed := failed + 1
      say("FAIL " + label)
    end
  end

  close_to(label: String, expected, actual, tolerance: Real) do
    check(label + " (expected " + fixed(expected, 3) + ", got " + fixed(actual, 3) + ")",
      (expected - actual).abs <= tolerance)
  end

  summary() do
    say("---")
    say(passed.to_string + " passed, " + failed.to_string + " failed")
    if failed > 0 then
      exit(1)
    end
  end
end

function pt(x, y: Real): Point do
  result := create Point.make(x, y)
end
\end{newfile}

\code{close_to} is for numbers worked out with \code{Real}s, which are
rarely exactly equal to what we expect, because of rounding. It passes if
the two are within a given tolerance.

\section{The basics}

The first checks are of the helpers, the geography and the geometry. Each
compares a result with an answer we know independently: one degree of
latitude is $2\pi R / 360$; a 10-metre square has an area of 100; the
point 5 metres along a line that goes 3 metres east and then 4 north is
2 metres up its second segment.

\begin{append}{checks.nex}
function check_helpers(c: Checker) do
  c.check("fixed rounds", fixed(2.345, 2) = "2.35" and fixed(-0.5, 1) = "-0.5")
  c.check("distance_text", distance_text(347.0) = "350 m" and distance_text(1234.0) = "1.2 km")
  c.check("duration_text", duration_text(20.0) = "under a minute"
    and duration_text(150.0) = "3 min" and duration_text(3900.0) = "1 h 05 min")
end

function check_geo(c: Checker) do
  let a: Lat_Lon := create Lat_Lon.make(47.0, -34.0)
  let b: Lat_Lon := create Lat_Lon.make(48.0, -34.0)
  -- A degree of latitude is 1/360 of a circle around the Earth.
  c.close_to("one degree of latitude, in metres", 2.0 * Real.PI * Earth.RADIUS / 360.0,
    a.distance_to(b), 0.01)
  c.close_to("distance is symmetric", a.distance_to(b), b.distance_to(a), 0.000001)
  let p: Projection := create Projection.make(a)
  let q: Lat_Lon := create Lat_Lon.make(47.02, -33.97)
  c.close_to("projection keeps distances, a few km from the centre", a.distance_to(q),
    p.to_point(a).distance_to(p.to_point(q)), a.distance_to(q) * 0.001)
  let back: Lat_Lon := p.to_lat_lon(p.to_point(q))
  c.close_to("projection round trip", 0.0, back.distance_to(q), 0.001)
  c.check("east is +x, north is +y", p.to_point(q).x > 0.0 and p.to_point(q).y > 0.0)
end

function check_geometry(c: Checker) do
  let square: Array[Point] := [pt(0.0, 0.0), pt(10.0, 0.0), pt(10.0, 10.0), pt(0.0, 10.0)]
  c.close_to("area of a 10 m square", 100.0, polygon_area(square), 0.000001)
  c.check("inside the square", polygon_contains(square, pt(5.0, 5.0)))
  c.check("outside the square", not polygon_contains(square, pt(15.0, 5.0))
    and not polygon_contains(square, pt(-1.0, 5.0)))
  c.close_to("distance to a segment, beside it", 5.0,
    distance_to_segment(pt(5.0, 5.0), pt(0.0, 0.0), pt(10.0, 0.0)), 0.000001)
  c.close_to("distance to a segment, past its end", 5.0,
    distance_to_segment(pt(13.0, 4.0), pt(0.0, 0.0), pt(10.0, 0.0)), 0.000001)
  c.check("bearings", bearing(pt(0.0, 0.0), pt(0.0, 1.0)) = 0.0
    and bearing(pt(0.0, 0.0), pt(1.0, 0.0)) = 90.0
    and bearing(pt(0.0, 0.0), pt(0.0, -1.0)) = 180.0
    and bearing(pt(0.0, 0.0), pt(-1.0, 0.0)) = 270.0)
  c.check("turn angles wrap around north", turn_angle(350.0, 10.0) = 20.0
    and turn_angle(10.0, 350.0) = -20.0 and turn_angle(0.0, 180.0) = 180.0)
  c.check("compass names", compass_name(0.0) = "north" and compass_name(100.0) = "east"
    and compass_name(350.0) = "north")
  let line: Array[Point] := [pt(0.0, 0.0), pt(3.0, 0.0), pt(3.0, 4.0)]
  c.close_to("line length", 7.0, line_length(line), 0.000001)
  c.close_to("point along a line", 2.0, point_along(line, 5.0).y, 0.000001)
end
\end{append}

\section{The map and its graphs}

Next, the map. Every place should be on land, and not in the lake; the
middle of the lake should be water; a point out at sea should be in no
area at all. And the graphs should obey the rules of the road: Market
Street is one-way, eastbound, so the driving graph must not have an edge
going west along it, while the walking graph must; and no footpath may
appear in the driving graph.

\begin{append}{checks.nex}
function check_map(c: Checker, m: Island_Map) do
  c.check("map loaded", m.name = "Kestrel Island" and m.places.length = 10
    and m.roads.length > 0 and m.areas.length > 0)
  let all_on_land: Boolean := true
  across m.places as p do
    if ?m.area_at(m.point(p.node_id)) as a then
      if a.kind = "water" then
        all_on_land := false
      end
    else
      all_on_land := false
    end
  end
  c.check("every place is on land", all_on_land)
  let lake_centre: Point := pt(0.0, 0.0)
  across m.areas as a do
    if a.kind = "water" then
      lake_centre := average(m.points_of(a.node_ids))
    end
  end
  if ?m.area_at(lake_centre) as a then
    c.check("the middle of the lake is water", a.kind = "water")
  else
    c.check("the middle of the lake is water", false)
  end
  c.check("out at sea is no area", m.area_at(pt(m.max_x() + 100.0, 0.0)) = nil)
  c.check("find_place by part of a name", ?m.find_place("mill") as p and p.name = "Old Mill")
end

function check_graphs(c: Checker, m: Island_Map, driving, walking: Road_Graph) do
  c.check("walking can use more roads than driving", walking.edge_count > driving.edge_count)
  -- Market Street is one-way, eastbound.
  let eastward: Boolean := true
  let westward_on_foot: Boolean := false
  across m.roads as r do
    if r.name = "Market Street" then
      across driving.edges_from(r.node_ids.get(1)) as e do
        if e.road.name = "Market Street" and e.target = r.node_ids.get(0) then
          eastward := false
        end
      end
      across walking.edges_from(r.node_ids.get(1)) as e do
        if e.road.name = "Market Street" and e.target = r.node_ids.get(0) then
          westward_on_foot := true
        end
      end
    end
  end
  c.check("cars may not drive the wrong way up a one-way street", eastward)
  c.check("walkers may", westward_on_foot)
  let no_footpaths: Boolean := true
  across driving.nodes as id do
    across driving.edges_from(id) as e do
      if e.road.kind = Road_Kind.Footpath then
        no_footpaths := false
      end
    end
  end
  c.check("no footpaths in the driving graph", no_footpaths)
end
\end{append}

\section{The searches}

Now the important ones. For every pair of places, and for drivers and
walkers alike, we run every search, and check:

\begin{itemize}[nosep]
  \item \textbf{Every place can reach every other.} With one-way streets,
    it is easy to make a map where a car can get into a street but never
    out. A real map gets this wrong surprisingly often.
  \item \textbf{A* agrees with Dijkstra's algorithm}, for the quickest
    routes and for the shortest. This is what A*'s mathematics promises,
    and if our estimate were ever too high, it is where we would see it.
  \item \textbf{Breadth-first search finds the fewest steps}: no route the
    others find has fewer.
  \item \textbf{The estimates never overshoot}: the estimate from the
    start is never more than the true cost.
  \item \textbf{A* does less work}, over all the pairs together.
  \item \textbf{The triangle inequality}: for any three places $a$, $b$
    and $x$, the shortest way from $a$ to $b$ is no longer than the
    shortest way from $a$ to $x$ plus the shortest from $x$ to $b$. If it
    ever were, the way through $x$ would be a shorter way from $a$ to $b$
    that the search had missed.
\end{itemize}

That is 100 pairs of places (each place to itself included) and five
searches for each, twice over: a thousand searches. They take a few
seconds.

\begin{append}{checks.nex}
-- The node a traveller in graph g starts from to go to or from a place.
function node_of(g: Road_Graph, p: Place): Integer do
  result := g.nearest_node(g.island.point(p.node_id))
end

function check_searches(c: Checker, m: Island_Map, g: Road_Graph, label: String) do
  -- Shortest distances between every pair of places, for the triangle
  -- check below.
  let dist: Map[String, Real] := {}
  let all_found: Boolean := true
  let same_time: Boolean := true
  let same_distance: Boolean := true
  let bfs_fewest: Boolean := true
  let admissible: Boolean := true
  let a_star_work: Integer := 0
  let dijkstra_work: Integer := 0
  across m.places as a do
    across m.places as b do
      let s: Integer := node_of(g, a)
      let t: Integer := node_of(g, b)
      let fast: Search_Result := dijkstra(g, s, t, by_time())
      let fast_a: Search_Result := a_star(g, s, t, by_time(), time_estimate(g, t))
      let short: Search_Result := dijkstra(g, s, t, by_distance())
      let short_a: Search_Result := a_star(g, s, t, by_distance(), distance_estimate(g, t))
      let fewest: Search_Result := breadth_first(g, s, t)
      all_found := all_found and fast.found() and short.found() and fewest.found()
      same_time := same_time and (fast.cost - fast_a.cost).abs < 0.001
      same_distance := same_distance and (short.cost - short_a.cost).abs < 0.001
      bfs_fewest := bfs_fewest and fewest.path.length <= short.path.length
        and fewest.path.length <= fast.path.length
      admissible := admissible and time_estimate(g, t)(s) <= fast.cost + 0.001
        and distance_estimate(g, t)(s) <= short.cost + 0.001
      a_star_work := a_star_work + fast_a.explored.length
      dijkstra_work := dijkstra_work + fast.explored.length
      dist.put(a.name + ">" + b.name, short.cost)
    end
  end
  c.check(label + ": every place can reach every other", all_found)
  c.check(label + ": A* finds routes as quick as Dijkstra's", same_time)
  c.check(label + ": A* finds routes as short as Dijkstra's", same_distance)
  c.check(label + ": breadth-first search finds the fewest steps", bfs_fewest)
  c.check(label + ": the estimates never overshoot", admissible)
  c.check(label + ": A* explores fewer nodes than Dijkstra ("
    + a_star_work.to_string + " against " + dijkstra_work.to_string + ")",
    a_star_work < dijkstra_work)
  let triangle: Boolean := true
  across m.places as a do
    across m.places as b do
      across m.places as x do
        triangle := triangle and dist.get(a.name + ">" + b.name)
          <= dist.get(a.name + ">" + x.name) + dist.get(x.name + ">" + b.name) + 0.001
      end
    end
  end
  c.check(label + ": no detour through a third place is shorter", triangle)
end
\end{append}

\code{time_estimate(g, t)(s)} makes the estimate function for target
\code{t} and calls it straight away, on \code{s}.

The shortest distances are remembered in a \code{Map} with keys like
\code{"Old Mill>Lighthouse"}, so the triangle check, with its thousand
combinations, needs no more searches.

Last, the program itself:

\begin{append}{checks.nex}
function main() do
  let c: Checker := create Checker
  let m: Island_Map := create Island_Map.load(user_path("kestrel_island.map"))
  let driving: Road_Graph := create Road_Graph.build(m, Mode.Driving)
  let walking: Road_Graph := create Road_Graph.build(m, Mode.Walking)
  check_helpers(c)
  check_geo(c)
  check_geometry(c)
  check_map(c, m)
  check_graphs(c, m, driving, walking)
  check_searches(c, m, driving, "driving")
  check_searches(c, m, walking, "walking")
  c.summary()
end

main()
\end{append}

\runoutput{checks.nex}{checks_core}

Everything passes. Look at the work done: over all the pairs, A* explores
about half as many nodes as Dijkstra's algorithm for drivers, and about a
third for walkers, for whom the straight-line estimate is closer to the
truth, just as the last chapter said.

Run the checks again whenever you change anything: the map file, a speed,
an algorithm. If you break something, a line saying \code{FAIL} tells you
what.

Try it. In \code{time_estimate}, divide by \code{speed / 2.0} instead of
\code{speed}, so that the estimate is twice what it should be, and run the
checks again. They don't even get as far as the search checks:

\begin{textout}
Error: Postcondition violation: cost_adds_up
  a_star, called ... in check_searches, does not deliver its postcondition
  `cost_adds_up`.
\end{textout}

With an estimate that is too high, A* can finish a node before it has found
the cheapest way to it. When the cheaper way turns up later, the search
gives the node a new parent and a new cost, but the nodes it already
reached from it keep the costs they got through the old way. The path
traced back through the parents then no longer costs what the search
says it does, and \code{a_star}'s own postcondition catches the
contradiction the first time it happens, long before anyone could notice a
slightly wrong route on the map. Put \code{speed} back, and everything
passes again.
